\section{Methodology}

\subsection{Motivation: Offline Variance Profiling}

Our approach is motivated by the GRPO gradient starvation identity: within-group reward standard deviation $\sigma = \sqrt{k(G-k)}/G$ equals zero when $k=0$ or $k=G$. For binary execution rewards, the expected variance for problem $i$ with pass rate $p_i$ is:

\begin{equation}
\mathrm{Var}_i = p_i \cdot (1 - p_i)
\end{equation}

This is maximized at $p_i = 0.5$ and equals zero for trivially easy or trivially hard problems. If we can identify problems near $p_i \approx 0.5$ before training, we can concentrate gradient steps on the productive fraction without paying the per-step cost of online difficulty estimation. For short RLEF regimes (20--50 steps), frozen-model performance should be a stable enough proxy --- the model's capabilities do not change dramatically over such a short window.

\subsection{Offline Variance Profiling Protocol}

\begin{enumerate}
\item Load the frozen base model (DeepSeek-Coder-7B-Instruct, weights frozen).
\item For each problem $i$ in MBPP training split (374 problems), generate $k=4$ independent completions using vLLM batch inference (temperature=1.0, top\_p=0.95, \texttt{max\_new\_tokens=128}).
\item Execute each completion against test cases via subprocess execution (5-second timeout). Record binary outcome: 1 if all tests pass, 0 otherwise.
\item Compute pass rate $p_i = (\text{correct completions}) / k$.
\item Compute variance estimate: $v_i = p_i \cdot (1 - p_i)$.
\end{enumerate}

\textbf{Implementation:} vLLM batch inference (\texttt{gpu\_memory\_utilization=0.75}) reduces profiling from $\sim$8 hours (sequential HF at $k=8$, \texttt{max\_new\_tokens=512}) to $\sim$32 seconds for $374\times4=1{,}496$ completions on H100 NVL.

\textbf{Parameter note:} The design specified $k=8$ and \texttt{max\_new\_tokens=512}. Hardware constraints and vLLM-TRL incompatibility (vLLM 0.11.0 + TRL 1.10) required reduction to $k=4$, \texttt{max\_new\_tokens=128} for profiling while using \texttt{max\_completion\_length=512} in training. This parameter mismatch is a central finding (Section~\ref{sec:results}).

\subsection{Variance-Guided Selection}

Select the top-$k$ problems by variance:

\begin{equation}
S_{\mathrm{variance}} = \mathrm{argsort}(v_i, \text{descending})[:k]
\end{equation}

where $k=50$. This selects problems where the model's pass rate is closest to 0.5 --- the zone of maximum gradient signal.

\textbf{Baseline:} random-50 selection --- 50 problems drawn uniformly at random (seed=42). This is the natural null hypothesis and the current default for practitioners without profiling infrastructure.

\textbf{Boundary behavior:} With $k=4$ discrete pass rates under \texttt{max\_new\_tokens=128}, all intermediate-success problems cluster at variance=0.1875. When fewer than 50 problems have variance $> 0$, the top-50 includes zero-variance problems at the boundary (\texttt{boundary\_gap=0}).

\subsection{GRPO Training Protocol}

We train with TRL GRPOTrainer (version 1.9.2):

\begin{table}[h]
\centering
\caption{GRPO training hyperparameters.}
\label{tab:training}
\begin{tabular}{ll}
\toprule
Parameter & Value \\
\midrule
Model & \texttt{deepseek-ai/deepseek-coder-7b-instruct-v1.5} \\
num\_generations ($G$) & 4 \\
max\_completion\_length & 512 \\
learning\_rate & $5\times10^{-7}$ (h-m2); $1\times10^{-6}$ (h-m3) \\
max\_steps & 50 (h-m2); 200 (h-m3) \\
beta (KL penalty) & 0.0 \\
use\_vllm & False \\
\bottomrule
\end{tabular}
\end{table}

\textbf{Reward function:} Binary execution reward via subprocess execution (5-second timeout): 1.0 if all test cases pass, 0.0 otherwise.

\textbf{Key diagnostic:} TRL automatically logs \texttt{frac\_reward\_zero\_std} --- the fraction of training groups where all $G=4$ completions have identical reward, producing $\sigma=0$. We monitor this metric as the primary diagnostic for whether data selection is having any effect.
